\documentclass{article}
\usepackage{enumitem}
\usepackage{amsmath}
\begin{document}

\title{Chapter 24: Sexual Reproduction in Flowering Plants}
\date{}
\maketitle

\begin{enumerate}[label=Q\arabic*.]

\item Double fertilization in angiosperms involves:
\\(A) Two pollen grains fertilizing one egg
\\(B) One sperm fusing with egg (forming zygote) and another sperm fusing with two polar nuclei (forming triploid endosperm)
\\(C) Fertilization occurring twice in the same flower
\\(D) Two eggs being fertilized by two sperm separately

\item The megaspore mother cell (MMC) undergoes meiosis to produce:
\\(A) Four functional megaspores
\\(B) One functional megaspore and three degenerate megaspores (in most angiosperms)
\\(C) Eight nuclei directly
\\(D) Two functional megaspores

\item The mature female gametophyte (embryo sac) of angiosperms typically contains how many cells/nuclei?
\\(A) 4 cells / 4 nuclei
\\(B) 7 cells / 8 nuclei (egg, 2 synergids, 3 antipodals, central cell with 2 polar nuclei)
\\(C) 8 cells / 8 nuclei
\\(D) 4 cells / 8 nuclei

\item Pollen kitt is a sticky substance that:
\\(A) Nourishes pollen during development
\\(B) Coats pollen grains and helps them adhere to pollinators and stigma
\\(C) Inhibits self-pollination
\\(D) Provides nutrients to the embryo sac

\item Assertion: Xenogamy ensures maximum genetic recombination.\\
Reason: Xenogamy involves transfer of pollen from one plant to the stigma of a genetically different plant of the same species.
\\(A) Both true, Reason is correct explanation
\\(B) Both true, Reason is not correct explanation
\\(C) Assertion true, Reason false
\\(D) Assertion false, Reason true

\item Cleistogamous flowers ensure:
\\(A) Cross-pollination
\\(B) Obligate self-pollination as flowers never open
\\(C) Wind pollination
\\(D) Insect pollination in closed buds

\item The pericarp of a fruit develops from:
\\(A) Ovule wall
\\(B) Ovary wall
\\(C) Thalamus
\\(D) Petals

\item Apomixis is significant in agriculture because:
\\(A) It produces genetically variable seeds
\\(B) It allows production of seeds without fertilization, maintaining elite genotypes
\\(C) It increases cross-pollination efficiency
\\(D) It enables vegetative propagation

\item Polyembryony in citrus occurs due to:
\\(A) Multiple fertilizations
\\(B) Development of additional embryos from nucellus cells (nucellar polyembryony) besides the zygotic embryo
\\(C) Apomixis only
\\(D) Fragmentation of zygote

\item In angiosperms, the pollen tube enters the ovule through which structure in porogamy?
\\(A) Chalaza
\\(B) Micropyle
\\(C) Integuments
\\(D) Funicle

\item Self-incompatibility (SI) in plants prevents self-fertilization by:
\\(A) Producing pollen that cannot germinate on any stigma
\\(B) S-gene mediated recognition and rejection of self-pollen by stigma
\\(C) Temporal separation of anther dehiscence and stigma receptivity (dichogamy)
\\(D) Physical separation of stamens and carpel

\item The endosperm in angiosperms is:
\\(A) Haploid nutritive tissue
\\(B) Triploid nutritive tissue providing nutrition to developing embryo
\\(C) Diploid embryonic tissue
\\(D) Formed from the egg cell after double fertilization

\item Which of the following is an agent of anemophily (wind pollination)?
\\(A) \textit{Salvia} (by bees)
\\(B) \textit{Vallisneria} (by water)
\\(C) \textit{Maize} (corn)
\\(D) \textit{Ficus} (by wasps)

\item Match the following parts of a seed with their origin:\\
1. Seed coat (testa) $\rightarrow$ (a) Integuments of ovule\\
2. Endosperm $\rightarrow$ (b) Triple fusion product (triploid)\\
3. Embryo $\rightarrow$ (c) Fertilized egg (zygote)\\
4. Cotyledons $\rightarrow$ (d) Part of embryo derived from zygote
\\(A) 1-a, 2-b, 3-c, 4-d
\\(B) 1-b, 2-a, 3-d, 4-c
\\(C) 1-c, 2-d, 3-a, 4-b
\\(D) 1-d, 2-c, 3-b, 4-a

\item The phenomenon of parthenocarpy results in:
\\(A) Seed development without fertilization
\\(B) Fruit development without fertilization, producing seedless fruits
\\(C) Multiple embryos in one seed
\\(D) Development of pollen grains without meiosis

\end{enumerate}
\end{document}
